\documentclass[aspectratio=169, 11pt]{beamer}

% ─────────────────────────────────────────────────────────
% THÈME ET PACKAGES
% ─────────────────────────────────────────────────────────
\usetheme{Madrid}
\usecolortheme{seahorse}
\setbeamertemplate{navigation symbols}{}
\setbeamertemplate{footline}[frame number]

\usepackage[utf8]{inputenc}
\usepackage[T1]{fontenc}
\usepackage[french]{babel}
\usepackage{amsmath, amssymb}
\usepackage{listings}
\usepackage{xcolor}
\usepackage{tikz}
\usetikzlibrary{arrows.meta, positioning, calc, decorations.pathreplacing, matrix}
\usepackage{booktabs}
\usepackage{graphicx}
\usepackage{hyperref}

% ─────────────────────────────────────────────────────────
% STYLE CODE PYTHON
% ─────────────────────────────────────────────────────────
\definecolor{codebg}{RGB}{245,245,250}
\definecolor{codegreen}{RGB}{0,128,0}
\definecolor{codepurple}{RGB}{128,0,128}
\definecolor{codered}{RGB}{200,30,30}
\definecolor{accentblue}{RGB}{0,80,160}
\definecolor{accentorange}{RGB}{230,120,0}

\lstdefinestyle{python}{
    language=Python,
    basicstyle=\ttfamily\scriptsize,
    backgroundcolor=\color{codebg},
    keywordstyle=\color{accentblue}\bfseries,
    stringstyle=\color{codered},
    commentstyle=\color{codegreen}\itshape,
    numberstyle=\tiny\color{gray},
    numbers=left,
    numbersep=5pt,
    frame=single,
    rulecolor=\color{gray!30},
    breaklines=true,
    showstringspaces=false,
    tabsize=4,
    morekeywords={self, torch, nn, math, None, True, False},
    emph={ScaledDotProductAttention, MultiHeadAttention, forward, __init__},
    emphstyle=\color{codepurple}\bfseries,
}

% ─────────────────────────────────────────────────────────
% MÉTADONNÉES
% ─────────────────────────────────────────────────────────
\title[Attention \& Multi-Head]{Le Mécanisme d'Attention\\et Multi-Head Attention}
\subtitle{De l'intuition au code --- Comprendre le cœur du Transformer}
\author{Master 1 LCS --- Lecture et Communication Scientifique}
\institute{D'après \emph{An Introduction to Transformers} (Turner, 2024)}
\date{Février 2026}

\begin{document}

% ═══════════════════════════════════════════════════════════
% PAGE DE TITRE
% ═══════════════════════════════════════════════════════════
\begin{frame}
\titlepage
\end{frame}

% ═══════════════════════════════════════════════════════════
% PLAN
% ═══════════════════════════════════════════════════════════
\begin{frame}{Plan de la présentation}
\tableofcontents
\end{frame}

% ═══════════════════════════════════════════════════════════
% PARTIE 1 : RAPPEL & MOTIVATION
% ═══════════════════════════════════════════════════════════
\section{Rappel : du signal aux patches}

\begin{frame}{Rappel : ce qu'on a vu}
\begin{columns}
\begin{column}{0.55\textwidth}
    \textbf{Étape précédente : le Patching}
    \begin{itemize}
        \item Signal de 100 points $\rightarrow$ 20 patches de 5
        \item Chaque patch est projeté en dimension $d_{\text{model}} = 64$
        \item On y ajoute un \emph{positional encoding}
    \end{itemize}
    
    \vspace{0.5em}
    \textbf{Résultat :} une matrice $X \in \mathbb{R}^{20 \times 64}$
    
    \vspace{0.5em}
    \alert{Question :} Comment chaque patch ``sait-il'' ce que contiennent les autres ?
\end{column}
\begin{column}{0.42\textwidth}
    \begin{tikzpicture}[scale=0.65, every node/.style={font=\footnotesize}]
        % Signal
        \draw[thick, blue] (0,0) -- (1,0) -- (1.5,2) -- (2,0) -- (3,0) -- (3.2,1.5) -- (3.5,1.5) -- (3.7,0) -- (5,0);
        \node[above, blue] at (2.5, -0.6) {Signal d'entrée};
        
        % Flèche
        \draw[->, thick] (2.5,-1) -- (2.5,-1.8);
        
        % Patches
        \foreach \i/\c in {0/red!60, 1/orange!60, 2/green!60, 3/blue!60, 4/purple!60} {
            \fill[\c, rounded corners=1pt] (\i*1, -2.2) rectangle (\i*1+0.85, -2.8);
        }
        \node at (2.5, -3.2) {5 patches (sur 20)};
        
        % Flèche
        \draw[->, thick] (2.5,-3.6) -- (2.5,-4.4);
        
        % Matrice
        \draw[thick] (0.5,-4.6) rectangle (4.5,-5.8);
        \node at (2.5, -5.2) {$X \in \mathbb{R}^{20 \times 64}$};
        
        % Flèche vers attention
        \draw[->, thick, red!70!black] (2.5,-6.1) -- (2.5,-6.7);
        \node[red!70!black, font=\bfseries] at (2.5,-7.1) {ATTENTION ?};
    \end{tikzpicture}
\end{column}
\end{columns}
\end{frame}

\begin{frame}{Pourquoi l'attention ?}
\begin{block}{Le problème}
Après le patching, chaque patch est traité \textbf{indépendamment}. \\
Or pour moyenner les pics par type, un patch-triangle doit ``savoir'' la hauteur des \emph{autres} patches-triangles.
\end{block}

\vspace{0.5em}

\begin{alertblock}{La solution : l'attention}
Un mécanisme qui permet à chaque position de \textbf{consulter toutes les autres} positions et d'en extraire l'information pertinente.
\end{alertblock}

\vspace{0.5em}

\begin{exampleblock}{Analogie (Alammar, 2018)}
\textbf{``it''} dans \emph{``The animal didn't cross the street because \textbf{it} was too tired''}  \\[0.3em]
$\rightarrow$ L'attention permet au modèle d'associer ``it'' à ``animal'' plutôt qu'à ``street''.
\end{exampleblock}
\end{frame}

% ═══════════════════════════════════════════════════════════
% PARTIE 2 : SCALED DOT-PRODUCT ATTENTION
% ═══════════════════════════════════════════════════════════
\section{Scaled Dot-Product Attention}

\begin{frame}{L'intuition : Query, Key, Value}
\begin{columns}
\begin{column}{0.55\textwidth}
    \textbf{Analogie de la bibliothèque} (Alammar) :
    
    \begin{enumerate}
        \item \textbf{Query} (Q) : \emph{``Je cherche un livre sur X''}
        \item \textbf{Key} (K) : \emph{L'étiquette de chaque livre}
        \item \textbf{Value} (V) : \emph{Le contenu de chaque livre}
    \end{enumerate}
    
    \vspace{0.5em}
    
    \textbf{Processus :}
    \begin{itemize}
        \item On compare la \textbf{query} avec chaque \textbf{key} $\rightarrow$ score de pertinence
        \item On pondère chaque \textbf{value} par ce score
        \item On fait la somme pondérée $\rightarrow$ résultat
    \end{itemize}
\end{column}
\begin{column}{0.42\textwidth}
    \begin{tikzpicture}[
        scale=0.55,
        every node/.style={font=\scriptsize},
        qbox/.style={fill=red!20, rounded corners, minimum width=1.5cm, minimum height=0.6cm},
        kbox/.style={fill=blue!20, rounded corners, minimum width=1.5cm, minimum height=0.6cm},
        vbox/.style={fill=green!20, rounded corners, minimum width=1.5cm, minimum height=0.6cm},
    ]
        % Query
        \node[qbox] (q) at (0, 4) {Query $q_1$};
        
        % Keys
        \node[kbox] (k1) at (3.5, 5.5) {Key $k_1$};
        \node[kbox] (k2) at (3.5, 4) {Key $k_2$};
        \node[kbox] (k3) at (3.5, 2.5) {Key $k_3$};
        
        % Scores
        \draw[->, thick] (q) -- (k1) node[midway, above, sloped] {0.7};
        \draw[->, thick] (q) -- (k2) node[midway, above] {0.2};
        \draw[->, thick] (q) -- (k3) node[midway, below, sloped] {0.1};
        
        % Values
        \node[vbox] (v1) at (7, 5.5) {Val $v_1$};
        \node[vbox] (v2) at (7, 4) {Val $v_2$};
        \node[vbox] (v3) at (7, 2.5) {Val $v_3$};
        
        \draw[->, thick, green!50!black] (k1) -- (v1);
        \draw[->, thick, green!50!black] (k2) -- (v2);
        \draw[->, thick, green!50!black] (k3) -- (v3);
        
        % Résultat
        \node[fill=yellow!30, rounded corners, minimum width=2cm, minimum height=0.6cm] (out) at (7, 0.8) {$0.7v_1 + 0.2v_2 + 0.1v_3$};
        
        \draw[->, thick, orange] (v1) -- (out);
        \draw[->, thick, orange] (v2) -- (out);
        \draw[->, thick, orange] (v3) -- (out);
    \end{tikzpicture}
\end{column}
\end{columns}
\end{frame}

\begin{frame}{La formule mathématique}
\begin{block}{Scaled Dot-Product Attention (Vaswani et al., 2017)}
\vspace{-0.5em}
\[
\boxed{\text{Attention}(Q, K, V) = \text{softmax}\!\left(\frac{QK^\top}{\sqrt{d_k}}\right) V}
\]
\end{block}

\vspace{0.3em}

\textbf{Décomposition étape par étape :}
\begin{enumerate}
    \item \textbf{Scores} : $S = QK^\top$ \hfill produit scalaire entre chaque query et chaque key
    \item \textbf{Scaling} : $S' = S / \sqrt{d_k}$ \hfill empêche les gradients de s'effondrer
    \item \textbf{Softmax} : $W = \text{softmax}(S')$ \hfill normalise en probabilités (somme = 1)
    \item \textbf{Pondération} : $\text{output} = W \cdot V$ \hfill somme pondérée des values
\end{enumerate}

\vspace{0.5em}

\begin{columns}
\begin{column}{0.48\textwidth}
\textbf{Dimensions :}\\
$Q \in \mathbb{R}^{n \times d_k}$, \;
$K \in \mathbb{R}^{m \times d_k}$, \;
$V \in \mathbb{R}^{m \times d_v}$\\
$\text{output} \in \mathbb{R}^{n \times d_v}$
\end{column}
\begin{column}{0.48\textwidth}
\textbf{Pourquoi $\sqrt{d_k}$ ?} (Turner, §3.2)\\
Si $q, k \sim \mathcal{N}(0,1)$, alors $q \cdot k$ a variance $d_k$.\\
$\rightarrow$ Diviser par $\sqrt{d_k}$ ramène la variance à 1.
\end{column}
\end{columns}
\end{frame}

\begin{frame}[fragile]{Le code : \texttt{ScaledDotProductAttention}}
\begin{lstlisting}[style=python, title={\texttt{src/transformers/attention.py} --- Lignes 56--97}]
class ScaledDotProductAttention(nn.Module):
    def __init__(self, dropout: float = 0.0):
        super().__init__()
        self.dropout = nn.Dropout(p=dropout)

    def forward(self, query, key, value, mask=None):
        d_k = query.size(-1)  # dimension des keys

        # Etape 1+2: scores = Q @ K^T / sqrt(d_k)
        scores = torch.matmul(query, key.transpose(-2, -1)) \
                 / math.sqrt(d_k)

        # Masquage optionnel (padding ou causal)
        if mask is not None:
            scores = scores.masked_fill(mask, float('-inf'))

        # Etape 3: softmax -> poids d'attention
        attention_weights = F.softmax(scores, dim=-1)
        attention_weights = self.dropout(attention_weights)

        # Etape 4: somme ponderee des values
        output = torch.matmul(attention_weights, value)
        return output, attention_weights
\end{lstlisting}
\end{frame}

\begin{frame}{Visualisation : que fait l'attention ?}
\begin{columns}
\begin{column}{0.5\textwidth}
    \begin{tikzpicture}[scale=0.5, every node/.style={font=\tiny}]
        % Matrice Q*K^T
        \def\data{
            {0.9, 0.1, 0.0, 0.0, 0.0},
            {0.1, 0.4, 0.1, 0.3, 0.1},
            {0.0, 0.1, 0.8, 0.0, 0.1},
            {0.0, 0.3, 0.0, 0.4, 0.3},
            {0.0, 0.1, 0.1, 0.3, 0.5}
        }
        \foreach \row [count=\y] in \data {
            \foreach \val [count=\x] in \row {
                \pgfmathsetmacro{\intensity}{\val*100}
                \fill[blue!\intensity!white] (\x-1, 5-\y) rectangle (\x, 5-\y+1);
                \node at (\x-0.5, 5-\y+0.5) {\pgfmathprintnumber[fixed, precision=1]{\val}};
            }
        }
        \draw[thick] (0,0) rectangle (5,5);
        
        % Labels
        \node[rotate=90] at (-0.5, 2.5) {Query (patch)};
        \node at (2.5, -0.4) {Key (patch)};
        \node[font=\small\bfseries] at (2.5, 5.6) {Poids d'attention (softmax)};
        
        % Patch labels
        \foreach \i in {1,...,5} {
            \node at (\i-0.5, -0.8) {$p_\i$};
            \node at (-1, 5.5-\i) {$p_\i$};
        }
    \end{tikzpicture}
\end{column}
\begin{column}{0.48\textwidth}
    \textbf{Lecture de la matrice :}
    
    \begin{itemize}
        \item Chaque \textbf{ligne} = comment un patch répartit son attention
        \item Ligne 1 : le patch 1 s'attend surtout à \textbf{lui-même} (0.9)
        \item Ligne 2 : le patch 2 regarde les patches 2 \emph{et} 4 (0.4 + 0.3)
    \end{itemize}
    
    \vspace{0.5em}
    \textbf{Pour notre tâche toy\_seq2seq :}
    \begin{itemize}
        \item Les patches-triangles devraient s'``écouter'' entre eux
        \item Les patches-carrés de même
        \item $\rightarrow$ Pour calculer la moyenne par type !
    \end{itemize}
\end{column}
\end{columns}
\end{frame}

\begin{frame}{Pourquoi $\sqrt{d_k}$ est crucial --- Démonstration}
\begin{columns}
\begin{column}{0.55\textwidth}
    \textbf{Sans scaling} ($d_k = 64$) :
    \begin{itemize}
        \item $q \cdot k = \sum_{i=1}^{64} q_i k_i$
        \item Si $q_i, k_i \sim \mathcal{N}(0,1)$ : variance = $d_k = 64$
        \item Scores $\approx \pm 16$ (grand en magnitude)
        \item Softmax $\rightarrow$ quasi one-hot : \textbf{gradients $\approx$ 0}
    \end{itemize}
    
    \vspace{0.5em}
    \textbf{Avec scaling} :
    \begin{itemize}
        \item $\frac{q \cdot k}{\sqrt{64}} = \frac{q \cdot k}{8}$
        \item Variance ramenée à $\approx 1$
        \item Scores $\approx \pm 2$ (raisonnable)
        \item Softmax $\rightarrow$ distribution douce : \textbf{gradients sains}
    \end{itemize}
\end{column}
\begin{column}{0.42\textwidth}
    \begin{tikzpicture}[scale=0.6, every node/.style={font=\scriptsize}]
        % Axe
        \draw[->] (0,0) -- (6,0) node[right] {$x$};
        \draw[->] (0,0) -- (0,4.5) node[above] {softmax};
        
        % Softmax douce (avec scaling)
        \draw[thick, blue, smooth] plot coordinates {(0.5,0.3) (1.5,0.5) (2.5,1.2) (3,2.5) (3.5,3.5) (4,3.8) (4.5,3.9) (5.5,4)};
        \node[blue, right] at (5.2, 3.5) {avec $\sqrt{d_k}$};
        
        % Softmax dure (sans scaling)
        \draw[thick, red, smooth] plot coordinates {(0.5,0) (1.5,0) (2,0) (2.5,0.05) (2.8,0.5) (3,3.5) (3.2,3.95) (3.5,4) (5.5,4)};
        \node[red, right] at (5.2, 4.3) {sans scaling};
        
        % Annotation
        \draw[<->, thick, orange] (2.8, 1.5) -- (3.2, 1.5);
        \node[orange, below] at (3, 1.3) {gradient};
    \end{tikzpicture}
    
    \vspace{0.3em}
    \centering
    \small\textcolor{gray}{Sans scaling : la softmax sature,\\les gradients disparaissent.}
\end{column}
\end{columns}
\end{frame}

% ═══════════════════════════════════════════════════════════
% PARTIE 3 : MULTI-HEAD ATTENTION
% ═══════════════════════════════════════════════════════════
\section{Multi-Head Attention}

\begin{frame}{Pourquoi plusieurs têtes ?}
\begin{block}{Limitation d'une seule tête}
Avec une seule attention, le résultat est une \textbf{moyenne pondérée unique}. \\
Le modèle ne peut capturer qu'\textbf{un seul type de relation} à la fois.
\end{block}

\vspace{0.5em}

\begin{exampleblock}{Multi-Head = plusieurs ``points de vue'' simultanés}
Chaque tête apprend des \textbf{projections différentes} $(W^Q_i, W^K_i, W^V_i)$ \\
$\rightarrow$ chacune se spécialise sur un aspect différent :
\end{exampleblock}

\vspace{0.3em}

\begin{columns}
\begin{column}{0.45\textwidth}
    \textbf{Tête 1 :} Relations de position\\
    \emph{``Qui est mon voisin ?''}
    
    \vspace{0.3em}
    \textbf{Tête 2 :} Similarité de forme\\
    \emph{``Qui a la même forme que moi ?''}
\end{column}
\begin{column}{0.45\textwidth}
    \textbf{Tête 3 :} Comparaison de hauteur\\
    \emph{``Qui a une hauteur similaire ?''}
    
    \vspace{0.3em}
    \textbf{Tête 4 :} Vision globale\\
    \emph{``Quel est le contexte général ?''}
\end{column}
\end{columns}

\vspace{0.5em}

\centering
\small \textcolor{gray}{(Alammar : ``it'' $\rightarrow$ tête 1 regarde ``animal'', tête 2 regarde ``tired'')}
\end{frame}

\begin{frame}{La formule Multi-Head}
\begin{block}{Multi-Head Attention (Vaswani et al., 2017)}
\vspace{-0.5em}
\[
\boxed{\text{MultiHead}(Q, K, V) = \text{Concat}(\text{head}_1, \dots, \text{head}_h) \, W^O}
\]
\[
\text{où} \quad \text{head}_i = \text{Attention}(Q W_i^Q, \; K W_i^K, \; V W_i^V)
\]
\end{block}

\vspace{0.3em}

\textbf{Dimensions avec $h$ têtes} ($h = 4$ dans notre code, $h = 8$ dans le papier) :

\vspace{0.3em}

\begin{center}
\begin{tabular}{lll}
\toprule
\textbf{Matrice} & \textbf{Dimension} & \textbf{Rôle} \\
\midrule
$W_i^Q, W_i^K, W_i^V$ & $\mathbb{R}^{d_{\text{model}} \times d_k}$ & Projections par tête \\
$d_k = d_{\text{model}} / h$ & $64 / 4 = 16$ & Dimension par tête \\
$W^O$ & $\mathbb{R}^{h \cdot d_k \times d_{\text{model}}}$ & Recombinaison \\
\bottomrule
\end{tabular}
\end{center}

\vspace{0.3em}
\textbf{Coût :} identique à une seule attention en $d_{\text{model}}$ ! \\
$h$ têtes de dimension $d_k$ = 1 tête de dimension $d_{\text{model}}$ (même nombre de multiplications).
\end{frame}

\begin{frame}{Pipeline visuel : Multi-Head Attention}
\centering
\begin{tikzpicture}[
    scale=0.55, 
    every node/.style={font=\scriptsize},
    box/.style={draw, rounded corners, minimum width=1.8cm, minimum height=0.7cm, thick},
    arrow/.style={->, thick, >=Stealth},
]
    % Input
    \node[box, fill=gray!20] (input) at (0, 0) {$X \in \mathbb{R}^{20 \times 64}$};
    
    % Projections
    \node[box, fill=red!20] (wq) at (-4, -2) {$W^Q$};
    \node[box, fill=blue!20] (wk) at (0, -2) {$W^K$};
    \node[box, fill=green!20] (wv) at (4, -2) {$W^V$};
    
    \draw[arrow] (input) -- (wq);
    \draw[arrow] (input) -- (wk);
    \draw[arrow] (input) -- (wv);
    
    % Split
    \node[font=\scriptsize\itshape, gray] at (0, -3.2) {split en $h=4$ têtes};
    
    % Heads
    \foreach \i/\col in {1/red!40, 2/orange!40, 3/green!40, 4/blue!40} {
        \pgfmathsetmacro{\xpos}{-4.5 + (\i-1)*3}
        \node[box, fill=\col] (head\i) at (\xpos, -4.5) {Tête \i};
        \node[font=\tiny, below=0.1cm of head\i] {$\mathbb{R}^{20 \times 16}$};
    }
    
    \draw[arrow] (wq) -- (head1);
    \draw[arrow] (wq) -- (head2);
    \draw[arrow] (wk) -- (head3);
    \draw[arrow] (wv) -- (head4);
    
    % Attention par tête
    \foreach \i/\col in {1/red!40, 2/orange!40, 3/green!40, 4/blue!40} {
        \pgfmathsetmacro{\xpos}{-4.5 + (\i-1)*3}
        \node[box, fill=\col!70] (attn\i) at (\xpos, -6.5) {\tiny Attention};
    }
    
    \foreach \i in {1,...,4} {
        \draw[arrow] (head\i) -- (attn\i);
    }
    
    % Concat
    \node[box, fill=yellow!30, minimum width=5cm] (concat) at (0, -8.2) {Concat $\rightarrow$ $\mathbb{R}^{20 \times 64}$};
    
    \foreach \i in {1,...,4} {
        \draw[arrow] (attn\i) -- (concat);
    }
    
    % W^O
    \node[box, fill=purple!20] (wo) at (0, -9.7) {$W^O$ : projection finale};
    \draw[arrow] (concat) -- (wo);
    
    % Output
    \node[box, fill=gray!20] (output) at (0, -11.2) {Sortie $\in \mathbb{R}^{20 \times 64}$};
    \draw[arrow] (wo) -- (output);
\end{tikzpicture}
\end{frame}

\begin{frame}[fragile]{Le code : \texttt{MultiHeadAttention} (1/2) --- Init}
\begin{lstlisting}[style=python, title={\texttt{src/transformers/attention.py} --- Projections Q, K, V}]
class MultiHeadAttention(nn.Module):
    def __init__(self, d_model: int, num_heads: int, dropout=0.0):
        super().__init__()
        assert d_model % num_heads == 0  # divisibilite obligatoire

        self.d_model = d_model        # 64
        self.num_heads = num_heads     # 4
        self.d_k = d_model // num_heads  # 64/4 = 16 par tete

        # 4 matrices de projection lineaire
        self.W_q = nn.Linear(d_model, d_model)  # 64 -> 64
        self.W_k = nn.Linear(d_model, d_model)  # 64 -> 64
        self.W_v = nn.Linear(d_model, d_model)  # 64 -> 64
        self.W_o = nn.Linear(d_model, d_model)  # 64 -> 64

        # Le mecanisme d'attention de base
        self.attention = ScaledDotProductAttention(dropout)
\end{lstlisting}

\vspace{0.3em}
\textbf{Astuce d'implémentation :} On utilise \texttt{nn.Linear(64, 64)} unique plutôt que 4 $\times$ \texttt{nn.Linear(64, 16)}. Le split en têtes se fait par \texttt{reshape} !
\end{frame}

\begin{frame}[fragile]{Le code : \texttt{MultiHeadAttention} (2/2) --- Forward}
\begin{lstlisting}[style=python, title={\texttt{src/transformers/attention.py} --- Forward pass}]
def forward(self, query, key, value, mask=None):
    batch_size = query.size(0)

    # 1. Projections lineaires
    Q = self.W_q(query)  # (batch, 20, 64)
    K = self.W_k(key)
    V = self.W_v(value)

    # 2. Reshape: split d_model en num_heads x d_k
    #    (batch, 20, 64) -> (batch, 4, 20, 16)
    Q = Q.view(batch_size, -1, self.num_heads, self.d_k) \
         .transpose(1, 2)
    K = K.view(batch_size, -1, self.num_heads, self.d_k) \
         .transpose(1, 2)
    V = V.view(batch_size, -1, self.num_heads, self.d_k) \
         .transpose(1, 2)

    # 3. Attention en parallele sur les 4 tetes
    attn_output, weights = self.attention(Q, K, V, mask)

    # 4. Concat: (batch, 4, 20, 16) -> (batch, 20, 64)
    attn_output = attn_output.transpose(1, 2).contiguous() \
                             .view(batch_size, -1, self.d_model)

    # 5. Projection finale W^O
    return self.W_o(attn_output), weights
\end{lstlisting}
\end{frame}

% ═══════════════════════════════════════════════════════════
% PARTIE 4 : LE RESHAPE EXPLIQUÉ
% ═══════════════════════════════════════════════════════════
\section{L'astuce du reshape}

\begin{frame}{Le reshape décortiqué}
\begin{center}
\begin{tikzpicture}[
    scale=0.6, 
    every node/.style={font=\scriptsize},
    mat/.style={draw, thick, minimum width=0.6cm, minimum height=0.4cm},
]
    % Étape 1: Après projection linéaire
    \node[font=\small\bfseries] at (0, 5) {Après $W^Q$ :};
    \draw[thick, fill=red!10] (-3, 3.5) rectangle (3, 4.5);
    \node at (0, 4) {\texttt{(batch, 20, 64)}};
    \node[gray] at (0, 3.2) {20 patches, chacun en 64 dims};
    
    % Flèche
    \draw[->, thick] (0, 2.8) -- (0, 2) node[midway, right] {\texttt{.view(..., 4, 16)}};
    
    % Étape 2: View
    \draw[thick, fill=orange!10] (-3, 0.8) rectangle (3, 1.8);
    \node at (0, 1.3) {\texttt{(batch, 20, 4, 16)}};
    \node[gray] at (0, 0.5) {20 patches $\times$ 4 têtes $\times$ 16 dims};
    
    % Flèche
    \draw[->, thick] (0, 0.1) -- (0, -0.7) node[midway, right] {\texttt{.transpose(1, 2)}};
    
    % Étape 3: Transpose
    \draw[thick, fill=green!10] (-3, -2) rectangle (3, -1);
    \node at (0, -1.5) {\texttt{(batch, 4, 20, 16)}};
    \node[gray] at (0, -2.5) {4 têtes $\times$ 20 patches $\times$ 16 dims};
    
    % Explication
    \node[align=left, font=\small] at (8, 2) {
        \textbf{Pourquoi ?}\\[0.3em]
        La dimension ``têtes'' doit\\
        être \emph{avant} la séquence\\
        pour que \texttt{torch.matmul}\\
        fasse l'attention en\\
        parallèle sur les 4 têtes.\\[0.3em]
        \textbf{C'est du batched matmul !}
    };
    
\end{tikzpicture}
\end{center}
\end{frame}

% ═══════════════════════════════════════════════════════════
% PARTIE 5 : SELF-ATTENTION, CROSS-ATTENTION, MASQUAGE
% ═══════════════════════════════════════════════════════════
\section{Variantes : Self, Cross, et Masquage}

\begin{frame}{Trois usages de l'attention dans le Transformer}
\begin{columns}
\begin{column}{0.32\textwidth}
    \begin{block}{\centering Self-Attention\\(Encoder)}
    \centering
    $Q = K = V = X$
    
    \vspace{0.3em}
    \footnotesize
    Chaque position regarde \textbf{toutes} les positions.
    
    \vspace{0.3em}
    \textit{C'est ce qu'utilise notre} \texttt{toy\_seq2seq} !
    \end{block}
\end{column}
\begin{column}{0.32\textwidth}
    \begin{block}{\centering Masked Self-Attn\\(Decoder)}
    \centering
    $Q = K = V = Y$
    
    \vspace{0.3em}
    \footnotesize
    Position $i$ ne voit que les positions $\leq i$.
    
    \vspace{0.3em}
    \textit{Masque causal (triangulaire supérieur $\rightarrow -\infty$).}
    \end{block}
\end{column}
\begin{column}{0.32\textwidth}
    \begin{block}{\centering Cross-Attention\\(Decoder)}
    \centering
    $Q = Y$, $K = V = \text{enc}$
    
    \vspace{0.3em}
    \footnotesize
    Le decoder \textbf{interroge} l'encoder.
    
    \vspace{0.3em}
    \textit{Les keys/values viennent de la source.}
    \end{block}
\end{column}
\end{columns}

\vspace{0.8em}

\begin{center}
\begin{tikzpicture}[every node/.style={font=\scriptsize}]
    % Masque causal
    \foreach \i in {0,...,4} {
        \foreach \j in {0,...,4} {
            \pgfmathsetmacro{\ok}{\j <= \i ? 1 : 0}
            \ifnum\ok=1
                \fill[green!30] (\j*0.5, -\i*0.5) rectangle (\j*0.5+0.5, -\i*0.5+0.5);
            \else
                \fill[red!30] (\j*0.5, -\i*0.5) rectangle (\j*0.5+0.5, -\i*0.5+0.5);
                \node at (\j*0.5+0.25, -\i*0.5+0.25) {$-\infty$};
            \fi
        }
    }
    \draw[thick] (0, -2) rectangle (2.5, 0.5);
    \node[below] at (1.25, -2.3) {Masque causal du decoder};
    
    \node at (4, -0.75) {\footnotesize Vert = visible};
    \node at (4, -1.25) {\footnotesize Rouge = masqué};
\end{tikzpicture}
\end{center}
\end{frame}

% ═══════════════════════════════════════════════════════════
% PARTIE 6 : LIEN AVEC NOTRE PROJET
% ═══════════════════════════════════════════════════════════
\section{Application : notre toy\_seq2seq}

\begin{frame}[fragile]{Dans notre modèle : \texttt{SimpleSignalTransformer}}
\begin{lstlisting}[style=python, title={\texttt{src/examples/toy\_seq2seq.py} --- Le modèle utilise l'attention}]
# L'encoder utilise nn.TransformerEncoderLayer qui contient:
#   - MultiHeadAttention (self-attention entre patches)
#   - FeedForward
#   - LayerNorm + Residual connections

encoder_layer = nn.TransformerEncoderLayer(
    d_model=64,          # dimension des embeddings
    nhead=4,             # 4 tetes d'attention
    dim_feedforward=256, # FFN interne
    dropout=0.1,
    batch_first=True
)
self.encoder = nn.TransformerEncoder(
    encoder_layer, num_layers=3  # 3 couches empilees
)
\end{lstlisting}

\vspace{0.3em}
\textbf{Chaîne complète :} \\
Signal $\xrightarrow{\text{patches}}$ $X \in \mathbb{R}^{20 \times 64}$ $\xrightarrow{\text{self-attention}}$ $\xrightarrow{\text{FFN}}$ $\xrightarrow{\times 3 \text{ couches}}$ Sortie
\end{frame}

\begin{frame}{Ce que les têtes apprennent dans notre tâche}
\begin{columns}
\begin{column}{0.48\textwidth}
    \textbf{Objectif du modèle :}
    \begin{itemize}
        \item Identifier le \emph{type} de chaque pic
        \item Calculer la \emph{moyenne} des hauteurs par type
        \item Remplacer chaque pic par cette moyenne
    \end{itemize}
    
    \vspace{0.5em}
    \textbf{Ce que l'attention permet :}
    \begin{itemize}
        \item Tête A : ``trouver les autres triangles''
        \item Tête B : ``trouver les autres carrés''
        \item La somme pondérée des values $\approx$ \textbf{moyenne}
    \end{itemize}
    
    \vspace{0.5em}
    \alert{Résultat :} 98.2\% de réduction de la loss !
\end{column}
\begin{column}{0.48\textwidth}
    \centering
    \begin{tikzpicture}[scale=0.55, every node/.style={font=\tiny}]
        % Input signal
        \node[font=\small\bfseries] at (3.5, 6.2) {Signal d'entrée};
        \draw[thick, blue] (0,5) -- (1,5) -- (1.3,5.8) -- (1.5,6) -- (1.7,5.8) -- (2,5) -- (3,5) -- (3.2,5.5) -- (3.5,5.5) -- (3.7,5) -- (4.5,5) -- (4.7,5.3) -- (4.9,5.5) -- (5.1,5.3) -- (5.3,5) -- (6,5) -- (6.2,5.8) -- (6.5,5.8) -- (6.7,5) -- (7,5);
        
        % Annotations
        \node[above, red] at (1.5, 6) {$\triangle$ h=10};
        \node[above, green!50!black] at (3.4, 5.5) {$\square$ h=5};
        \node[above, red] at (5, 5.5) {$\triangle$ h=6};
        \node[above, green!50!black] at (6.4, 5.8) {$\square$ h=8};
        
        % Attention arrows
        \draw[->, thick, red, dashed] (1.5, 4.8) to[bend right=30] (5, 4.8);
        \node[red, below] at (3.2, 4.3) {tête: ``même type''};
        
        \draw[->, thick, green!50!black, dashed] (3.4, 4.5) to[bend right=20] (6.4, 4.5);
        \node[green!50!black, below] at (4.9, 3.8) {tête: ``même type''};
        
        % Output
        \node[font=\small\bfseries] at (3.5, 2.8) {Sortie attendue};
        \draw[thick, orange] (0,1.5) -- (1,1.5) -- (1.3,2) -- (1.5,2.2) -- (1.7,2) -- (2,1.5) -- (3,1.5) -- (3.2,2.15) -- (3.5,2.15) -- (3.7,1.5) -- (4.5,1.5) -- (4.7,1.8) -- (4.9,2.2) -- (5.1,1.8) -- (5.3,1.5) -- (6,1.5) -- (6.2,2.15) -- (6.5,2.15) -- (6.7,1.5) -- (7,1.5);
        
        \node[above, red] at (1.5, 2.2) {avg=8};
        \node[above, green!50!black] at (3.4, 2.15) {avg=6.5};
        \node[above, red] at (5, 2.2) {avg=8};
        \node[above, green!50!black] at (6.4, 2.15) {avg=6.5};
    \end{tikzpicture}
\end{column}
\end{columns}
\end{frame}

% ═══════════════════════════════════════════════════════════
% PARTIE 7 : SYNTHÈSE
% ═══════════════════════════════════════════════════════════
\section{Synthèse et points clés}

\begin{frame}{Récapitulatif : ce qu'il faut retenir}
\begin{enumerate}
    \item \textbf{Attention = recherche d'information}
    \begin{itemize}
        \item Query interroge, Key répond, Value fournit le contenu
        \item Formule : $\text{softmax}(QK^\top / \sqrt{d_k}) \, V$
    \end{itemize}
    
    \vspace{0.3em}
    \item \textbf{Le scaling $\sqrt{d_k}$ est essentiel}
    \begin{itemize}
        \item Empêche la saturation du softmax
        \item Sans lui : gradients $\approx 0$, pas d'apprentissage
    \end{itemize}
    
    \vspace{0.3em}
    \item \textbf{Multi-Head = plusieurs perspectives}
    \begin{itemize}
        \item Coût identique grâce au split $d_k = d_{\text{model}} / h$
        \item Chaque tête capture un type de relation différent
    \end{itemize}
    
    \vspace{0.3em}
    \item \textbf{3 variantes} : self-attention, masked self-attention, cross-attention
    
    \vspace{0.3em}
    \item \textbf{Le reshape est l'astuce clé} de l'implémentation
    \begin{itemize}
        \item Pas $h$ boucles $\rightarrow$ un seul \texttt{matmul} batché
    \end{itemize}
\end{enumerate}
\end{frame}

\begin{frame}{Complexité et limites}
\begin{columns}
\begin{column}{0.55\textwidth}
    \textbf{Complexité} (Turner, §3 ; Fleuret) :
    \[
    \mathcal{O}(n^2 \cdot d)
    \]
    \begin{itemize}
        \item $n$ = nombre de tokens (patches)
        \item $d$ = dimension du modèle
        \item \textbf{Quadratique} en la longueur de séquence !
    \end{itemize}
    
    \vspace{0.5em}
    \textbf{C'est pourquoi le patching aide :}
    \begin{itemize}
        \item Sans : $n = 100 \Rightarrow 100^2 = 10\,000$
        \item Avec : $n = 20 \Rightarrow 20^2 = 400$
        \item $\rightarrow$ 25$\times$ moins de calcul !
    \end{itemize}
\end{column}
\begin{column}{0.42\textwidth}
    \textbf{Pistes modernes :}
    \begin{itemize}
        \item \emph{Linear attention} : $\mathcal{O}(n \cdot d^2)$
        \item \emph{Flash Attention} : optimise la mémoire
        \item \emph{Sparse attention} : ne regarde qu'un sous-ensemble
    \end{itemize}
    
    \vspace{0.5em}
    \textcolor{gray}{(Mentionné dans Turner, §5)}
\end{column}
\end{columns}
\end{frame}

% ═══════════════════════════════════════════════════════════
% SLIDE FINALE
% ═══════════════════════════════════════════════════════════

\begin{frame}{Références}
\footnotesize
\begin{thebibliography}{9}
    \bibitem{turner2024} Turner, R. E. (2024). \emph{An Introduction to Transformers}. arXiv:2304.10557v5.
    \bibitem{vaswani2017} Vaswani, A. et al. (2017). \emph{Attention Is All You Need}. NeurIPS.
    \bibitem{alammar2018} Alammar, J. (2018). \emph{The Illustrated Transformer}. \url{https://jalammar.github.io/illustrated-transformer/}
    \bibitem{rush2018} Rush, A. et al. (2018). \emph{The Annotated Transformer}. \url{https://nlp.seas.harvard.edu/2018/04/03/attention.html}
    \bibitem{fleuret2022} Fleuret, F. (2022). \emph{Transformers --- Slides}. \url{https://fleuret.org/public/EN_20220809-Transformers/}
\end{thebibliography}

\vspace{1em}

\centering
\textbf{Code source :} \texttt{src/transformers/attention.py}\\
\textbf{Exemple :} \texttt{python -m src.examples.toy\_seq2seq}
\end{frame}

\end{document}
